% Please do not change %
\documentclass[11pt]{article} % font size
\usepackage[margin=1in]{geometry} % margins
\usepackage{amsmath,amsthm,amssymb} % for the  common "math symbols"
\usepackage{algorithm}
\usepackage{algpseudocode}
\usepackage[shortlabels]{enumitem} % for enumeration
\usepackage{booktabs} % if you need tables
\usepackage{listings}
\usepackage{color}
\usepackage{quiver}
\usepackage{ulem}
\DeclareMathOperator*{\argmax}{argmax}
\DeclareMathOperator*{\argmin}{argmin}

% \theoremstyle{definition}
\newtheorem{theorem}{Theorem}[section]
\newtheorem{corollary}[theorem]{Corollary}
\newtheorem{lemma}[theorem]{Lemma}
\newtheorem{proposition}[theorem]{Proposition}
\newtheorem{conjecture}[theorem]{Conjecture}
\theoremstyle{definition}
\newtheorem{definition}[theorem]{Definition}
\theoremstyle{definition}
\newtheorem{fact}[theorem]{Fact}
\theoremstyle{definition}
\newtheorem{example}[theorem]{Example}
\newtheorem{note}[theorem]{Note}

\lstset{
  frame=none,
  xleftmargin=2pt,
  stepnumber=1,
  numbers=left,
  numbersep=5pt,
  numberstyle=\ttfamily\tiny\color[gray]{0.3},
  belowcaptionskip=\bigskipamount,
  captionpos=b,
  escapeinside={*'}{'*},
  language=haskell,
  tabsize=2,
  emphstyle={\bf},
  commentstyle=\it,
  stringstyle=\mdseries\rmfamily,
  showspaces=false,
  keywordstyle=\bfseries\rmfamily,
  columns=flexible,
  basicstyle=\small\sffamily,
  showstringspaces=false,
  morecomment=[l]\%,
}
% Feel free to include additional packages (if needed), but make sure it does not override the margin/font requirements.
%
\setlength{\parindent}{0pt} % no indent for new paragraphs.
\setlength{\parskip}{5pt} % distance between paragraphs, which will be used when you have a blank line in your LaTeX code.


%%%%%%%%%%%%%%%%%%%%%%%%%%%%%%%%%%%%%%%%%%%%%%%%%%%%%%%%%%%%%%%%%%%%%%%%%%%%%%
\title{Categories, Proofs, and Processes: \\Lecture 11 Naturality}
%
\author{Wira Azmoon Ahmad}
\date{8 Feb 2024}
%
\begin{document}
%
\maketitle

%You may use the usual \LaTeX\ environments to structure your answers:
%\begin{align}
%        \label{basegnn}
%        \mathbf{h}_u^{(t+1)} = \sigma \bigg( \mathbf{W}_\text{self}^{(t)} \mathbf{h}_u^{(t)}  + \mathbf{W}_\text{neigh}^{(t)} \sum_{v \in N(u)} \mathbf{h}_v^{(t)} + \mathbf{b}^{(t)} \bigg): 
%\end{align}


%\begin{table}[h]
%\caption{A simple table.}
%\centering
%\begin{tabular}[t]{lrrrr}
%\toprule
%Graphs & Can & Be  & Fun & Too \\ 
%\midrule 
%$G_1$ & $1$  & $2$ & $3$ & $4$ \\ 
%$G_2$ & $-1$  & $-2$ & $-3$ & $-4$ \\ 
%$G_3$ & $0.1$  & $0.2$ & $0.3$ & $0.4$ \\ 
%\bottomrule
%\end{tabular}
%\label{tab}
%\end{table}

% A functor $F:\mathcal{C} \rightarrow \mathcal{D}$ is
% \begin{itemize}
%     \item \textbf{faithful} if each map $F_{A,B}:\mathcal{C}(A,B) \rightarrow \mathcal{D}(F\ A, F\ B)$ is injective;
%     \item \textbf{full} if each map $F_{A,B}:\mathcal{C}(A,B) \rightarrow \mathcal{D}(F\ A, F\ B)$ is surjective
% \end{itemize}


\addtocounter{section}{3}
\section{Naturality}
Functors treat categories as objects. We can treat functors themselves as objects, so what are the arrows?

\subsection{Motivating Examples}
Suppose a category $\mathcal{C}$ has products, we have two isomorphic ternary products.
\[(A\times B)\times C \cong A\times (B\times C)\]
Note that this isomorphism does not really depend on the specific objects and arrows involved,
i.e. the following commutes

\subsubsection[]{Associators}
% https://q.uiver.app/#q=WzAsNCxbMCwwLCIoQVxcdGltZXMgQilcXHRpbWVzIEMiXSxbMSwwLCIoQSdcXHRpbWVzIEIpXFx0aW1lcyBDIl0sWzAsMSwiQVxcdGltZXMgKEJcXHRpbWVzIEMpIl0sWzEsMSwiQSdcXHRpbWVzIChCXFx0aW1lcyBDKSJdLFswLDFdLFswLDIsImkiLDJdLFsxLDMsImknIl0sWzIsM11d
\[\begin{tikzcd}
	{(A\times B)\times C} & {(A'\times B)\times C} \\
	{A\times (B\times C)} & {A'\times (B\times C)}
	\arrow[from=1-1, to=1-2]
	\arrow["i"', from=1-1, to=2-1]
	\arrow["{i'}", from=1-2, to=2-2]
	\arrow[from=2-1, to=2-2]
\end{tikzcd}\]
More generally we have
\[(-\times -)\times - : \mathcal{C} \times \mathcal{C} \times \mathcal{C} \rightarrow \mathcal{C}\]
\[-\times (-\times -) : \mathcal{C} \times \mathcal{C} \times \mathcal{C} \rightarrow \mathcal{C}\]

\subsubsection{Lists}
Consider the list constructor
\[x_0 : [x_1,x_2,...,x_n] = [x_0,x_1,x_2,...,x_n]\]
"Liste" constructor in reverse
\[[x_n,x_{n-1},...,x_1]:x_0 = [x_n,x_{n-1},...,x_1,x_0]\]

Consider the functor \texttt{List} : \textbf{Set} $\rightarrow$ \textbf{Set},
which maps a set $X$ to \texttt{List}$(X)$, the set of all finite lists with elements in $X$.
% https://q.uiver.app/#q=WzAsNCxbMCwwLCJcXHRleHR0dHtMaXN0fShYKSJdLFsxLDAsIlxcdGV4dHR0e0xpc3R9KFkpIl0sWzAsMSwiXFx0ZXh0dHR7TGlzdH0oWCkiXSxbMSwxLCJcXHRleHR0dHtMaXN0fShZKSJdLFsxLDMsIlxcdGV4dHR0e3JldmVyc2V9Il0sWzAsMSwiXFx0ZXh0dHR7TGlzdH0oZikiXSxbMCwyLCJcXHRleHR0dHtyZXZlcnNlfSIsMl0sWzIsMywiXFx0ZXh0dHR7TGlzdH0oZikiLDJdXQ==
\[\begin{tikzcd}
	{\texttt{List}(X)} & {\texttt{List}(Y)} \\
	{\texttt{List}(X)} & {\texttt{List}(Y)}
	\arrow["{\texttt{reverse}}", from=1-2, to=2-2]
	\arrow["{\texttt{List}(f)}", from=1-1, to=1-2]
	\arrow["{\texttt{reverse}}"', from=1-1, to=2-1]
	\arrow["{\texttt{List}(f)}"', from=2-1, to=2-2]
\end{tikzcd}\]

\subsubsection{More Examples}
Consider functiors $G:\underline{2}_{\rightrightarrows}\rightarrow \textbf{Set}$.
% https://q.uiver.app/#q=WzAsMyxbMSwwLCJcXGJ1bGxldCJdLFsyLDAsIlxcYnVsbGV0Il0sWzAsMCwiXFx1bmRlcmxpbmV7Mn1fe1xccmlnaHRyaWdodGFycm93c30iXSxbMCwxLCIiLDIseyJjdXJ2ZSI6LTF9XSxbMCwxLCIiLDAseyJjdXJ2ZSI6MX1dLFsyLDAsIj0iLDEseyJzdHlsZSI6eyJib2R5Ijp7Im5hbWUiOiJub25lIn0sImhlYWQiOnsibmFtZSI6Im5vbmUifX19XV0=
\[\begin{tikzcd}
	{\underline{2}_{\rightrightarrows}} & \bullet & \bullet
	\arrow[curve={height=-6pt}, from=1-2, to=1-3]
	\arrow[curve={height=6pt}, from=1-2, to=1-3]
	\arrow["{=}"{description}, draw=none, from=1-1, to=1-2]
\end{tikzcd}\] 
corresponds to directed graphs
% https://q.uiver.app/#q=WzAsMixbMCwwLCJFIl0sWzEsMCwiViJdLFswLDEsInMiLDAseyJjdXJ2ZSI6LTF9XSxbMCwxLCJ0IiwyLHsiY3VydmUiOjF9XV0=
\[\begin{tikzcd}
	E & V
	\arrow["s", curve={height=-6pt}, from=1-1, to=1-2]
	\arrow["t"', curve={height=6pt}, from=1-1, to=1-2]
\end{tikzcd}\]
where $E\subseteq V\times V$.
% https://q.uiver.app/#q=WzAsNixbMSwwLCJFIl0sWzIsMCwiViJdLFsxLDEsIkUnIl0sWzIsMSwiViciXSxbMCwwLCJcXHRleHR7R3JhcGggMX0iXSxbMCwxLCJcXHRleHR7R3JhcGggMn0iXSxbMCwxLCJzIl0sWzAsMiwiaCIsMl0sWzIsMywicyciLDJdLFsxLDMsImgiXV0=
\[\begin{tikzcd}
	{\text{Graph 1}} & E & V \\
	{\text{Graph 2}} & {E'} & {V'}
	\arrow["s", from=1-2, to=1-3]
	\arrow["h"', from=1-2, to=2-2]
	\arrow["{s'}"', from=2-2, to=2-3]
	\arrow["h", from=1-3, to=2-3]
\end{tikzcd}\]
is a graph homomorphism.
\subsubsection{Unnatural transformation}
\[\texttt{reverse}_{\texttt{Bool}}: \texttt{List}(X)\rightarrow \texttt{List}(X)
:= \begin{cases}
    \texttt{reverse list} & \text{if } X = \texttt{Bool}\\
    \texttt{id} & \text{otherwise}
\end{cases}\]
Consider the function $\texttt{odd}: \mathbb{N} \rightarrow \texttt{Bool}$
(True if $n\in\mathbb{N}$ is odd, False otherwise).
\[\texttt{reverse}_{\texttt{Bool}} \circ \texttt{list odd}([1,3,2]) = [F,T,T]\]
\[\texttt{list odd}\circ\texttt{reverse}_{\texttt{Bool}}  ([1,3,2]) = [T,T,F]\]
Because the function cares about the types inside, it's unnatural.
The diagram will not commute.

\clearpage
\subsection{Natural Transformations}
\begin{definition}
Let $F,G:\mathcal{C} \rightarrow \mathcal{D}$ be functors.
A \underline{natural transformation}
\[t:F\rightarrow G\]
is a family of morphisms in $\mathcal{D}$ indexed by objects $A$ of $\mathcal{C}$,
\[\{t_A:F\ A \rightarrow G\ A\}_{A\in \text{Ob}(\mathcal{C})}\]
such that, for all $f:A\rightarrow B$, the following diagram commutes
% https://q.uiver.app/#q=WzAsNCxbMCwwLCJGXFwgQSJdLFsxLDAsIkZcXCBCIl0sWzAsMSwiR1xcIEEiXSxbMSwxLCJHXFwgQiJdLFswLDEsIkZcXCBmIl0sWzAsMiwidF9BIiwyXSxbMiwzLCJHXFwgZiIsMl0sWzEsMywidF9CIl1d
\[\begin{tikzcd}
	{F\ A} & {F\ B} \\
	{G\ A} & {G\ B}
	\arrow["{F\ f}", from=1-1, to=1-2]
	\arrow["{t_A}"', from=1-1, to=2-1]
	\arrow["{G\ f}"', from=2-1, to=2-2]
	\arrow["{t_B}", from=1-2, to=2-2]
\end{tikzcd}\]
This condition is called naturality.
\end{definition}
\begin{definition}
    If each $t_A$ is an isomorphism, we say $t$ is a \underline{natural isomorphism}:
    \[t:F\cong G\]
\end{definition}

\begin{note}
	A natural isomorphism is precisely an isomorphism in the functor category 
	\[\text{Func}(\mathcal{C},\mathcal{D}):=\mathcal{D}^{\mathcal{C}}\]
\end{note}
	
\addtocounter{subsection}{3}
\subsection{Equivalence of categories}
\begin{definition}[Def. 57 (pg. 38)]
    An \underline{equivalence of categories} consists of a pair of functors
    \[E:\mathcal{C} \rightarrow \mathcal{D}\]
    \[F:\mathcal{D} \rightarrow \mathcal{C}\]
    and a pair of natural isomorphisms
    \[\alpha : 1_{\mathcal{C}} \cong F\circ E\]
    \[\beta : 1_{\mathcal{D}} \cong E\circ F\]
    $F$ is the \underline{pseudo-inverse} of $E$.
    $\mathcal{C},\mathcal{D}$ are then said to be equivalent, denoted $\mathcal{C}\simeq \mathcal{D}$.
\end{definition}
\begin{note}
    This is an \underline{alternative} definition to a functor being an equivalence.
\end{note}

\begin{proposition}[Ex. 1.4.4(2)(b)$\Rightarrow$(a)]
    If functor $F:\mathcal{C}\rightarrow\mathcal{D}$ is (part of) an equivalence of categories,
    it is an equivalence - full, faithful, and 
    \[\forall D\in\text{Ob}(\mathcal{D}),\exists C\in\text{Ob}(\mathcal{C}). F\ C \cong D\]
\end{proposition}
\begin{proof}
For faithfulness, take $E:\mathcal{D}\rightarrow \mathcal{C}$ such that
\[\alpha : 1_{\mathcal{C}} \cong EF\]
\[\beta : 1_{\mathcal{D}} \cong FE\]
are natural transformations in $\mathcal{C}$, for any object $C$, we have $\alpha_C : C \cong EF(C)$ and
% https://q.uiver.app/#q=WzAsNCxbMCwwLCJDIl0sWzEsMCwiQyciXSxbMCwxLCJFRihDKSJdLFsxLDEsIkVGKEMnKSJdLFswLDEsImYiXSxbMCwyLCJcXGFscGhhX0MiLDJdLFsyLDMsIkVGKGYpIiwyXSxbMSwzLCJcXGFscGhhX3tDJ30iXV0=
\[\begin{tikzcd}
	C & {C'} \\
	{EF(C)} & {EF(C')}
	\arrow["f", from=1-1, to=1-2]
	\arrow["{\alpha_C}"', from=1-1, to=2-1]
	\arrow["{EF(f)}"', from=2-1, to=2-2]
	\arrow["{\alpha_{C'}}", from=1-2, to=2-2]
\end{tikzcd}\]
Thus if $F(f)=F(f')$ then $EF(f)=EF(f')$ and so $f=f'$. So $F$ is faithful since
\[F_{C,C'}:\mathcal{C}(C,C')\rightarrow \mathcal{D}(F(C),F(C'))\]
is surjective.

For fullness, consider any arrow $h:F(C) \rightarrow F(C')$ in $\mathcal{D}$. 
We note that $C\cong EF(C)$ and $C'\cong EF(C')$. So consider $E(h)$ then there is an
\[f = (\alpha_{C'})^{-1} \circ E(h) \circ \alpha_C : C\rightarrow C'\]
where $\alpha_C,\alpha_{C'}$ are iso. By naturality we have $F(f):F(C)\rightarrow F(C')$ that
% https://q.uiver.app/#q=WzAsNCxbMCwwLCJDIl0sWzEsMCwiQyciXSxbMCwxLCJFRihDKSJdLFsxLDEsIkVGKEMnKSJdLFswLDEsImYiXSxbMCwyLCJcXGFscGhhX0MiLDJdLFsyLDMsIkVGKGYpIiwyXSxbMSwzLCJcXGFscGhhX3tDJ30iXV0=
\[\begin{tikzcd}
	C & {C'} \\
	{EF(C)} & {EF(C')}
	\arrow["f", from=1-1, to=1-2]
	\arrow["{\alpha_C}"', from=1-1, to=2-1]
	\arrow["{EF(f)}"', from=2-1, to=2-2]
	\arrow["{\alpha_{C'}}", from=1-2, to=2-2]
\end{tikzcd}\]
So $EF(f) = E(h) : EF(C) \rightarrow EF(C')$. Since $E$ is faithful, $EF(f)=E(h)\implies F(f) = h$, so $F$ is full.

For essential surjectivity on objects for any $D\in\text{Ob}(\mathcal{D})$, we have
\[t:1_{\mathcal{D}}\cong FE\]
so for any $D\in\text{Ob}(\mathcal{D})$,
\[\beta_D : D \cong FED = F(ED)\]
for $ED\in \text{Ob}(\mathcal{C})$.
\end{proof}

\subsection{Representable functors}
\begin{definition}
    A functor $F:\mathcal{C}\rightarrow\textbf{Set}$ is said to be \underline{representable}
    if it is naturally isomorphic to the covariant Hom-functor at $A$,
    $\mathcal{C}(A,-):\mathcal{C}\rightarrow\textbf{Set}$ for some $A$.
\end{definition}
\begin{example}
    Forgetful functor $U:\textbf{Mon}\rightarrow\textbf{Set}$.
    We want that there is a monoid $Z$ such that for all $M$,
    \[\textbf{Mon}(Z,M) \cong U(M), \quad Z(\mathbb{N}, +, 0)\]
\end{example}

\end{document}
